\section{Comparing two words in the selected spectral estimate}
\label{sec:offdiag-v62}

The distinct-word norm estimate uses both the level primes and the
archimedean column dimensions of the two words. The common endpoints
of the columns permit one constant $A=e^{3/4}$ to control both
ratios. Write $D=\mathcal F.D(x,i)$, $k=D.R.k$, $f=D.R.f$,
$\Delta=D.s.\mathrm{Del}$, $T_u=(D.\mathrm{mem}(u)).A.T$,
$L=x^\epsilon$ and $B=e^L$. All thresholds in $x$ below precede
the index $i$, the words and the places.

% BEGIN ODW1
\paragraph{ODW1. A logarithmic bound for the interior of a word.}
\label{claim:ODW1}
Fix $n\ge1$, $\Delta>0$ and $v_0,\ldots,v_{n-1}$ with
$v_{a-1}-v_a=4\Delta$ for $1\le a<n$. Given bits
$\xi_1,\ldots,\xi_n$, put $z_a=v_{a-1}-1/2-\Delta\xi_a$ and
$V(z)=\prod_{a<b}(z_a-z_b)$. For any two strings of bits,
\[
                    |\log V(z')-\log V(z)|\le 7n/12 .
\]
\begin{proof}
For $a<b$ divide the gap by $4\Delta(b-a)$. The resulting factor
is $1+t_{ab}$ with $t_{ab}=(\xi_b-\xi_a)/(4(b-a))$.
Integrating $t/(1+t)$ gives
$0\le t-\log(1+t)\le (2/3)t^2$ for $|t|\le1/4$.
Thus the nonlinear remainder in the log product belongs to
$[-n/12,0]$ for every word, since
\[
 \sum_{a<b}t_{ab}^2\le \frac{n}{16}\sum_{d\ge1}d^{-2}
 \le\frac{n}{8}.
\]
The last estimate follows by bounding the terms after the first by
$\int_1^\infty t^{-2}\,dt$. Two such remainders differ by at most
$n/12$.

Set $H_j=\sum_{d=1}^j d^{-1}$ and $H_0=0$. The linear sum is
$\tfrac14\sum_{a=1}^n\xi_a(H_{a-1}-H_{n-a})$.
The absolute coefficients before the factor $1/4$ have sum
\[
 2\sum_{a=1}^{\lfloor n/2\rfloor}\sum_{d=a}^{n-a}\frac1d
 \ \le\ 2\sum_{d=1}^{n-1}\frac{d}{d}\ \le\ 2n .
\]
For fixed $d$ there are at most $d$ admissible $a$.
As $|\xi_a'-\xi_a|\le1$, the two linear sums differ by at most
$n/2$. The common baseline product cancels. Adding the remainder
bound proves the assertion.
\end{proof}
% END ODW1

% BEGIN ODW2
\paragraph{ODW2. The full Vandermonde and the local critical levels.}
\label{claim:ODW2}
For every supplied $\mathrm{FamP}$, at all sufficiently large $x$,
every index, every pair of words and every place $\rho$ satisfy
\[
 \left|\log\frac{V_k(u'_{\rho})}{V_k(u_{\rho})}\right|
       \le \frac{3k}{4}-\frac{13}{12},\qquad
 \frac{T_{\rm crit}(u,\rho)}{T_{\rm crit}(u',\rho)}
       \le \exp\!\left(\frac34-\frac{13}{12k}\right).
\]
Here $u_{\rho}$ denotes the $k$ coordinates in $D.\mathrm{hcW}(u,\rho)$.
\begin{proof}
Put $n=k-2$. The identities of DSZ4 identify the interior with
ODW1, including its sign of the bit displacement. The fields
$D.\mathrm{hcW\_first}$ and $D.\mathrm{hcW\_last}$ fix the two
endpoints $U, V$ across all words at this place. We have
$U-v_0\ge4k\Delta$ and $V\le v_{k-2}-1/2$. Off the top places
these follow from the exact endpoint clauses of DSZ4. At a top
place $U\ge xB$, $0\le v_0\le4k\Delta\le 2B/x$, and $x\ge3$
suffices for the first inequality. Also
$V=-b_T-(k-1)/2\le-1/2$ since $b_T>0$, and $v_{k-2}\ge0$.
The endpoint-to-endpoint factor is common. For an upper endpoint
gap the smaller of its two values is at least $4k\Delta$, and
the change in the gap is at most $\Delta$. The total change of
upper endpoint log factors is therefore at most $n/(4k)\le1/4$.

For $1\le a\le n$ put $d=k-1-a\ge1$. At the lower endpoint
$z_a-V\ge(4d-1)\Delta$ because
$v_{a-1}-v_{k-2}=4d\Delta$. Hence the total lower endpoint change
is at most $\sum_{d=1}^n(4d-1)^{-1}\le H_n/3\le(n+1)/6$.
Both estimates use $|\log s'-\log s|\le |s'-s|/\min(s,s')$.
Adding the bounds gives
$7n/12+1/4+(n+1)/6=3k/4-13/12$.
All the factors are positive.

In the source formula for $T_{\rm crit}(u,\rho)^k$ the only
word-dependent factor is $1/V_k(u_\rho)$. The positive numerator is
the same host power times $D.\mathrm{dsig}(\rho)\prod_{j<k+m-1}j!$.
Taking the $k$th root gives the second inequality and, in
particular, the literal $\mathrm{cTvB}(D,e^{3/4})$.
\end{proof}
% END ODW2

% BEGIN ODW3
\paragraph{ODW3. Both word-size inequalities with one constant.}
\label{claim:ODW3}
Let $A=e^{3/4}$. Every supplied family eventually satisfies the
original $\mathrm{Sizes3b}(D,A)$, at all indices, with no
admissibility premise. In particular $a_F(l)=A$ for every list
$l$ supplies the whole closed $h_A$ input of the selected
distinct-word theorem on any class.
\begin{proof}
At $x>\mathcal F.x_1$ the extra record gives the following
identity for the positive integer stored with the word:
\[
 d_*(u)=(D.\mathrm{mem}(u)).A.\mathrm{duStar}
       =\mu_k I_{\rm col}\prod_\rho
       \frac{V_k(u_\rho)}{\prod_{j=0}^{k-1}j!}.
\]
Thus $d_*(u')/d_*(u)=\prod_\rho V_k(u'_\rho)/V_k(u_\rho)$,
with exactly $f$ factors. The center of the window for $T_u$ is
$S_u=(D_*/d_*(u))^{1/(fk)}\Gamma^{-2/(fk)}$.
The equalities $D.\mathrm{Dstar\_eq}$ and $D.\mathrm{GamP\_eq}$
give the common positive $D_*$ and $\Gamma$.
The window itself reads $S_u\le T_u\le(1+1/(fk))S_u$.
Consequently
\[
 \log(T_u/T_{u'})
 \le \frac1{fk}+\frac1{fk}\sum_\rho
                \log\frac{V_k(u'_\rho)}{V_k(u_\rho)}
 \le \frac34-\frac{13}{12k}+\frac1{fk}<\frac34 ,
\]
using $f\ge1$. This holds for both orderings.
ODW2 gives $V_k(u'_\rho)\le A^k V_k(u_\rho)$ at each place.
These are precisely the two clauses of $\mathrm{Sizes3b}$.
Taking the maximum of the DSZ4 threshold and
$\mathcal F.x_1+1$ proves $\mathrm{OwnEv}$.
For the closed assertion take $\epsilon_1=1$ and the zero
transported bound for every permitted $\epsilon$, before choosing
the family. The same event then holds at the family's own $C',C_l$
and under its source admissibility condition. No row-volume identity
or selection enters the prime ratio.
\end{proof}
% END ODW3

% BEGIN ODW4
\paragraph{ODW4. The counting hypotheses at the selected pair.}
\label{claim:ODW4}
Use $A$ from ODW3. At every sufficiently late index satisfying
admissibility with $C'\ge1$, $C_l>2$, the full
$\mathrm{Cnt.Hyp}$ holds for each ordered pair, including equal
words. In particular it applies to the two words of a selected
loop with $(C_g,c_b,C_u,c_T,C_T,C_0)=
(\mathcal F.\mathrm{CYL}+1,1/4,3,1/2,\mathcal F.\mathrm{CTu},3)$.
\begin{proof}
Admissibility supplies the standing numerical and size hypotheses
and $1\le C_0$. The even integer $m\ge C_l>2$ gives
$r_S=m/2-1\ge1$. Put $g=g_0g_1$. Both levels are primes,
split completely in $E$, and are $1$ modulo $g^2$. As $g^2>2$,
both primes are at least three. The host primes $g_0,g_1$ are
distinct and split, and each differs from the level primes. Thus
$\gcd(g,T_uT_{u'})=1$ and $T_u\equiv T_{u'}\pmod g$.
These are all clauses of $\mathrm{Cnt.Stand}$.

For the three remaining logarithmic estimates use $L>1$, $k\ge5$,
$m\ge2$, $m^4\le f$, $1\le C'$ and the original standing bounds:
\[
 C_lm^6L\le m^7L\le m^5fL\le C'm^5fL
       \le k\le k\log k .
\]
Also $0\le\log f\le\sqrt{r_S}$, so
$(\log f)^2\le r_S\le A m^2 L$ and $\log f\le A r_S^2$.
Together with ODW3 in both orders these prove each conjunct of
$\mathrm{Cnt.Hyp}$ for the same $T_u,T_{u'},D,x,i$. At class
scope, a constant final transported bound of $3$ enforces the two
stated size conditions on $sC',sC_l$.
\end{proof}
% END ODW4

% BEGIN ODW5
\paragraph{ODW5. A positive largest threshold at the actual places.}
\label{claim:ODW5}
For $C_A\ge1$ the closed $\mathrm{ThresholdAtPlaces}$ event holds on
any family class, for every admissible late index and every distinct
pair. Its threshold uses the actual least differing letter.
\begin{proof}
Use the canonical identification $e$ from infinite places of $E$ to
$D.\mathrm{what}_2.\mathrm{Tau}$, inverse to the map that takes the
source embedding to its infinite place. At each place use one-based
indices $v_j$, $1\le j\le k-1$, for the cloud. The letters of a word
are common across places. Admissibility makes the two columns different
at every place. Their endpoints agree, so there is a first differing
bit $l$, $1\le l\le k-2$. This is the source's $\mathrm{lOf}$.
Write $g=g_0g_1$, $\kappa=\pi/D.\mathrm{host}.v$, and write
$P(Y)=\prod_{s=1}^{k+m-1}(Y+\lambda_s)$ for the row polynomial at
this place. The denominator of the local function is
\[
 q(Y)=-T_uT_{u'}\prod_{a=1}^k
       (Y+1/2+u_a)(Y+1/2+u'_a).
\]
Consider the compact interval $5/2\le t=Y+v_l\le\Delta/2$.
It has positive length eventually since
$\Delta\ge N(m)m^{10}$ by DSZ4 and $m\longrightarrow\infty$.
On it each interior denominator factor has the form
$t+4\Delta(l-j)-\xi_j\Delta$. For $j<l$ it is positive and
for $j>l$ negative. For $j=l$ it is $t$ or $t-\Delta$, both
nonzero. At the upper endpoint of the column the factor is positive
by $U-v_1\ge4k\Delta$. At its lower endpoint the factor is at most
$t+v_{k-1}-v_l\le t-4\Delta<0$ by ODW2 and $l\le k-2$.
This proves $q(Y)\ne0$ throughout the interval.

The function
\[
  g^m\kappa^{-(m-1)}
     \sqrt{\left|P(Y)P(Y+1)/q(Y)\right|}
\]
is continuous there and is positive except at possibly finitely
many points, since $P$ is monic. Its supremum $M_w$ is thus
finite and positive at every infinite place $w$.
In the original $G.\mathrm{thrSet}$, take the
$\mathrm{lamPrimeHat}$ alternative at index $1$. Its value is
\[
             \prod_w\sqrt{C_A M_w\kappa/\sqrt{T_u T_{u'}}}>0 .
\]
The full set contains only finitely many real numbers: the primary
index runs from $1$ to $r_S+1$, with at most $r_S+2$ choices for
the secondary index in the branch that uses one. Let $M$ be its
maximum. Then $M>0$, $f>0$, and $\Theta=M^{-1/f}$ satisfies
$0<\Theta$ and $\Theta^{-f}=M$ with the required membership
and upper-bound properties. This argument only needs one positive
candidate in the full set. The interval lower bound is uniform
in all indices, places and pairs. The source's admissibility guard
can therefore be retained at the family's own constants, with zero
transported bound in the closed statement.
\end{proof}
% END ODW5

% BEGIN ODW6
\paragraph{ODW6. Masses, both ranks and the row normalization.}
\label{claim:ODW6}
On a family class with a supplied $\mathrm{TSel}$, the original
$\mathrm{SzRowsAt}$ holds eventually under the admissibility guard
at every index. At the same indices the numerical data in the
global size calculation satisfy the two column counts, the row
mass identity, the consecutive mass ratios and the $L$-value bound.
\begin{proof}
Write $s=r_S$, $d=k+m-1=k+2s+1$ and
$W_j=\prod_{a=0}^{j-1}a!$. The real and column records in the
selected $\mathrm{RecT}$ apply simultaneously, since their threshold
is common to all $i$. RSL3 applied to its real record identifies
the volume of $Y_0$ with the standard mass $\mu_d$.
RMP11 for its columns and the family's actual level record gives
the rank on the same $D$. With the extra record's stored $D_*$
these are the equalities
\[
 \begin{aligned}
 D.\mathrm{duS}(u)&=\mu_k I_{\rm col}
       \prod_\rho\frac{V_k(u_\rho)}{W_k}\,(T_u^k-1)^f,\\
 D_*&=\mu_d I_{\rm row}\prod_\rho D.\mathrm{dsig}(\rho).
 \end{aligned}
\]
The power $(T_u^k-1)^f$ multiplies the whole product. Both formulas
use the original indices. Indeed, complete splitting of $g_0,g_1$
in $E$ gives $f$ places of $F$ above each, all of norm $g_j$,
and the two sets of places are disjoint. Thus for any real-valued
$h$ the product of $h(\mathrm N\mathfrak p)$ over these typed
places is $h(g_0)^f h(g_1)^f$, exactly $\mathrm{SzaSpl}$.
For the local column index use $\prod_{t=1}^k(q^t-1)$.
For the row index divide $|\mathrm{GL}_d(\mathbb F_q)|$ by
$q^{s(2k+2s+1)+k(k-1)/2}$. Counting the free entries as in
QW2 proves these are respectively the local factors of
$\mathrm{gaG}$ and $\mathrm{idxRows}$. An arbitrary orthonormal
lattice in either stated rank is isometric to the standard one
by its given basis. Isometries preserve its genus and each stabilizer,
so its mass is $\mu_k$ or $\mu_d$. This verifies the universal
orthonormal-model quantifiers in $\mathrm{FamilyMidVolume}$.

At every rank $j\ge3$, MF19 on the unramified CM host gives the
prescribed ratio of orthonormal-model masses
\[
 \frac{\mu_j}{\mu_{j-1}}=
 \frac{((j-1)!(v/\pi)^j)^f}{\sqrt{(2v)^f}}L_E(j).
\]
This is $\mathrm{FamilySchurFormalDegree}$. DSZ5 on the same late
members gives $D.\mathrm{dsig}(\rho)W_d=V_d(\lambda_\rho)$.
In particular the required lower bound
$V_d(\lambda_\rho)/2\le D.\mathrm{dsig}(\rho)W_d$ holds at
every place. These four assertions constitute $\mathrm{SzRowsAt}$.

For the size data $G.\mathrm{size}$, reindex the finite products
along $e$. Its $\mathrm{vol}(j)$ is the mass of the orthonormal
telescope model, hence $\mu_j$. The preceding identities give its
$\mathrm{Dstar}$ and its $\mathrm{ColumnCount}$ for each of the
two words, with their own prime and their own dimension.
For $k<j\le d$ the displayed mass ratio has its designated
$d_F=(2v)^f$. QW1 gives $L_E(j)>0$ and
$|\log L_E(j)|\le f\log\zeta(j)$ for every such $j$.
Thus $\mathrm{MassRatios}$ and $\mathrm{LBound}$ hold in these
same size data. All uses of MF19 and DSZ5 carry their previously
recorded Hasse--Minkowski and formal-degree inputs.
\end{proof}
% END ODW6

% BEGIN ODW7
\paragraph{ODW7. The full selected size row.}
\label{claim:ODW7}
Let $\mathcal A$ have a supplied $\mathrm{TSel}$, and let
$\mathrm{sel}:\mathrm{SSel}(\mathcal A)$ be any supplied spectral
selector. The original $\mathrm{SizesRowS}(C_0,\mathcal A,\mathrm{sel})$
holds. Its selected $X,Y$, words, levels and dimensions are the ones
that enter the distinct-word bound.
\begin{proof}
Fix $C_A\ge1$ and arbitrary functions $k_0,c_n$ of the source's
constant list. Put $l_0=\mathcal F.cOut\mathbin{+\!+}
\mathcal F.cMid(C_H)$. In the closed event take $\epsilon_1=1$
and final transported bound $l\mapsto\max(256,8c_n(l))$, with
the other components zero. This puts both $C'=sC'(\mathcal F)$
and $C_l=sC_l(\mathcal F)$ above that maximum. Choose a common
$x$ threshold for ODW5--6, $B\ge16$, and $k\ge k_0(l_0)$
(the latter follows from $k\ge3x$ in DSZ4). This threshold may
depend on $\mathcal F,C_H$, and precedes all indices and words.

At an admissible point also satisfying the selector's two threshold
conditions choose any place $\tau_0$; there are $f>0$ places.
Use the same $e$ as in ODW5. The selected setting's defining
embedding identity $X.\phi_{\rm eq}$ identifies the embedding
$X.B.I.\phi(X.\mathrm{pl}^{-1}(ew))\circ X.\iota_E^{-1}$
with the source embedding for $ew$. Its infinite place is therefore
$w$. At the interior positions,
$(v_j-1/2-u_{j+1})/\Delta$ is exactly the bit of the word,
independent of $\tau_0$; outside them the source assigns zero.
Thus both word functions in $G.\mathrm{size}$, and their first
different position, agree with ODW5. Take the three scalar witnesses
$C_p=C'$, $\Theta$ of ODW5, and $C_f=256$.

We check the global shape conditions in $J.\mathrm{SzHyp}$.
The local configuration at $w$ uses exactly the actual row and
column endpoints at $ew$. The standing hypotheses give
$0\le v_{k-1}$, $v_1\le4k\Delta$, $U-v_1\ge4k\Delta$,
$1\le b_1$, and the top endpoint parameter
$b'=-V$ satisfies $b_1/2\le b'\le b_1-3$ at top places.
The bottom nodes have consecutive gaps at least one (DSZ4).
Together with $k\ge5$, $m\ge2$, $\Delta\ge1$ these are the
local $\mathrm{Ok}$ conditions. For local $\mathrm{Small}$,
$b_j\le B^2$, $4mk\Delta/b_1\le4/C_l\le1/10$, and
$2m\log B/k\le2/(C_lm)\le1/10$, using
$C_l m^2L\le k$. Globally $B>1$, $m\le m^4\le f\le k$,
$v\ge1$, $d_F=(2v)^f$, $(kB)^8\le g_j$, and the genus mass
$\mu_k$ is positive. The standing inequalities
$k^{C_l}\le b_j$ and $(kfL)^{C_l}\le\Delta$ give
$k^3\le b_j$ and $(kf\log B)^5\le\Delta$.
These verify $G.\mathrm{Hyp}(B,5,3)$.

Both bit functions take values $0,1$ and vanish at $k-1$,
and $1\le l\le k-2$ is their first differing position.
The standing assumptions give $B^{m/2}\le T_u,T_{u'}$,
$8C_\Gamma\le C_p$ and $C_pm^5 f\log B\le k$. ODW6 supplies
the mass ratios, $L$-value estimate, both column counts, and the
formal-degree lower bound. For each word the numerical standing
hypothesis gives the inequalities for
$D.\mathrm{duS}(u)$ between
$\Gamma^{-2}D_*/C_0$ and $C_0\Gamma^{-2}D_*$,
where the family field identifies
$\Gamma=\exp(C_\Gamma m^5 f\log B)$.
By ODW6 the model's $D_*$ is the same number. Thus these are
both $\mathrm{ChoiceOfT}$ clauses. ODW5 gives the positive
$\Theta$ and the greatest-element clause, completing
$J.\mathrm{SzHyp}$.

Finally $c_{\exp}(5,1/2,3)=1/16$ and
$256\log B\le C_l L\le f\log k$. Since $C'\ge256$,
$k\ge C'm^5 fL$ implies both $m+f+5\le k$ and
$4s+12=2m+8\le k$, where $s=r_S$. For
$\mathrm{PB.SizeNumbers}(k,s,f,B,g_0,g_1,1/3,c_n(l_0))$,
we already have $B\ge16$, $k\ge32$, $f\le k$, $4s+12\le k$,
and $(kB)^8\le g_j$. The same bound on $k$ gives
$70\log B\le k$ and $8c_n(l_0)m^5f\log B\le k$.
Also $1\le(5/3)f\log k$. The last two fields of
$\mathrm{SizesAt}$ are the identity for $D_*$ in ODW6 and
$T_u\equiv T_{u'}\equiv1\pmod{(g_0g_1)^2}$.
This supplies every witness for the source constants
$(C_A,C_0,\mathcal F.C_\Gamma,5,1/2,3,c_n(l_0),k_0(l_0))$,
for every admissible late $x,i$ and every selected distinct pair.
\end{proof}
% END ODW7

% BEGIN ODW8
\paragraph{ODW8. Remaining estimates in the selected norm inequality.}
\label{claim:ODW8}
On a supplied $\mathrm{TSel}$ class, fix the compatible spectral
selector of TFR3 and $C_0=3$. Assume the source's
$\mathrm{TheoremSSingle}$ and its selected
$\mathrm{LoopRowS}(3,\mathcal A,\mathrm{sel})$. Then the closed
$\mathrm{TheoremSSelected}$ holds on this class. In particular
the norm is that of the selected operator on the actual pair:
\[
 \|(Y.\mathrm{tO}(u,u')).K_0\|^2
  \le (1+B^{-6})\,
         \frac{D.\mathrm{duS}(u)D.\mathrm{duS}(u')}{D_*},
              \qquad u\ne u'.
\]
\begin{proof}
In $\mathrm{TheoremSSingle}$ use $C_{\rm thr}=1$, obtaining
$C_A\ge1$. At $(C_0,C_G,C_\Delta,c_T,C_b)=
(3,\mathcal F.C_\Gamma,5,1/2,3)$ the last product is $3/2>1$.
Its remaining constants are functions of the transported list,
since $\mathcal F.C_\Gamma,\mathcal F.\mathrm{CYL},
\mathcal F.\mathrm{CTu}$ occur at positions 17, 6 and 13
respectively (zero-based). Choose its $k_0,c_n$ and the lower
bounds for $C',C_l$ on these lists before the family. Increase
the last transported bound to include the two lower bounds and
29, and take its maximum with the bounds of ODW7 and of the
assumed loop assertion evaluated at $a_F\equiv e^{3/4}$.
Intersection of closed conditions takes the minimum of their
$\epsilon$ bounds and the componentwise maximum of their
transported bounds.

At a common late $x$ take the given admissibility guard and any
selected distinct pair. TFR2--3 and SFS6--7 give the gap,
$\mathrm{LocTied}$, $\mathrm{TFrameAt}$ and the concrete tie
for its very $X,Y$. QPR11 gives absolute convergence for that tie.
ODW4 gives $\mathrm{Cnt.Hyp}$ for the loop, and its
$\mathrm{Hyp0}$ and $\mathrm{Stand}$ fields for the pointwise
theorem. The accompanying $\mathrm{KTwoHypotheses}$ follow
from the same standing bounds: the split-prime bounds are
$(kB)^8\le g_j\le B^{C_g}$; summing positive adjacent row gaps
gives every required distinct-coordinate gap $>C_l k m^3$;
$C_l\log\mathrm{rd}_E\le L$ and $\log f\le r_S^2$ are the
stated or weaker logarithmic bounds. The level bounds and
coprimality come from admissibility and ODW4, with the unramified
CM condition from the host. The loop assumption supplies
$\mathrm{RowBig}$ and $\mathrm{RowNotScalar}$ at the specified
$\epsilon_L(D)$ and $\mathrm{ddC}(X,Y,u,u')$. ODW7 supplies
$\mathrm{SizesAt}$ for precisely this pair and $B>0$.
Substitution into $\mathrm{TheoremSSingle}$ gives the displayed
inequality. The maximum of the finitely many family thresholds
precedes $i,u,u'$ and works under both selector guards, as
required by $\mathrm{OwnEv}$. The two assumed estimates are the
outstanding pointwise spectral and selected row-loop inputs on
this route.
\end{proof}
% END ODW8
